\documentclass[11pt, a4paper]{article}

% Gemeinsame Style-Definitionen (TEXINPUTS muss templates/ enthalten — siehe scripts/build_tex.py)
% bfw-style.tex — gemeinsame Style-Definitionen für alle Handouts/Aufgabenhefte
% Voraussetzung: Kompilation mit xelatex (wegen fontspec)

\usepackage[ngerman]{babel}
\usepackage{geometry}
\geometry{a4paper, top=2.2cm, bottom=2.2cm, left=2.3cm, right=2.3cm}

\usepackage{fontspec}
% Fallback-Kette: Source Sans Pro -> Calibri -> Segoe UI -> Latin Modern Sans
% (Latin Modern ist immer mit MiKTeX vorhanden)
\IfFontExistsTF{Source Sans Pro}{%
  \setmainfont{Source Sans Pro}[Ligatures=TeX]%
}{\IfFontExistsTF{Calibri}{%
  \setmainfont{Calibri}[Ligatures=TeX]%
}{\IfFontExistsTF{Segoe UI}{%
  \setmainfont{Segoe UI}[Ligatures=TeX]%
}{%
  \setmainfont{Latin Modern Sans}[Ligatures=TeX]%
}}}
\IfFontExistsTF{Source Code Pro}{%
  \setmonofont{Source Code Pro}[Scale=0.92]%
}{\IfFontExistsTF{Consolas}{%
  \setmonofont{Consolas}[Scale=0.92]%
}{%
  \setmonofont{Latin Modern Mono}[Scale=0.92]%
}}

\usepackage{xcolor}
% Farbpalette: hell, freundlich, modern
\definecolor{bfwPrimary}{HTML}{0EA5A4}    % Petrol/Türkis - Primär
\definecolor{bfwPrimaryDark}{HTML}{0F766E} % dunkler Petrol für Headings
\definecolor{bfwAccent}{HTML}{F59E0B}     % Amber - Akzent
\definecolor{bfwSuccess}{HTML}{10B981}    % Grün - Erfolg / Lösung
\definecolor{bfwWarn}{HTML}{F97316}       % Orange - Achtung
\definecolor{bfwInk}{HTML}{0F172A}        % fast Schwarz
\definecolor{bfwInkSoft}{HTML}{475569}    % gedämpftes Grau
\definecolor{bfwBgSoft}{HTML}{F8FAFC}     % off-white
\definecolor{bfwBgTeal}{HTML}{ECFEFF}     % helles Türkis
\definecolor{bfwBgAmber}{HTML}{FEF3C7}    % helles Gelb
\definecolor{bfwBgRose}{HTML}{FFE4E6}     % helles Rosé für Eselsbrücken

\usepackage[colorlinks=true, linkcolor=bfwPrimaryDark, urlcolor=bfwPrimaryDark]{hyperref}
\usepackage{enumitem}
\setlist[itemize]{leftmargin=*, itemsep=2pt, topsep=2pt}
\setlist[enumerate]{leftmargin=*, itemsep=2pt, topsep=2pt}

\usepackage[most]{tcolorbox}
\usepackage{booktabs}
\usepackage{tikz}
\usetikzlibrary{decorations.pathmorphing, positioning, shapes.geometric, arrows.meta}

% Merkkästchen — wichtige Definition / Kernpunkt
\newtcolorbox{merksatz}[1][]{%
  enhanced, breakable,
  colback=bfwBgTeal,
  colframe=bfwPrimary,
  boxrule=0.6pt,
  arc=4pt,
  left=10pt, right=10pt, top=8pt, bottom=8pt,
  fonttitle=\bfseries\color{bfwPrimaryDark},
  title={\faLightbulb\ Merksatz},
  #1
}

% Eselsbrücke — Mnemoniks
\newtcolorbox{eselsbruecke}[1][]{%
  enhanced, breakable,
  colback=bfwBgRose,
  colframe=bfwAccent,
  boxrule=0.6pt,
  arc=4pt,
  left=10pt, right=10pt, top=8pt, bottom=8pt,
  fonttitle=\bfseries\color{bfwWarn},
  title={Eselsbrücke},
  #1
}

% Beispiel-Box
\newtcolorbox{beispiel}[1][]{%
  enhanced, breakable,
  colback=bfwBgSoft,
  colframe=bfwInkSoft,
  boxrule=0.4pt,
  arc=3pt,
  left=10pt, right=10pt, top=8pt, bottom=8pt,
  fonttitle=\bfseries\color{bfwInk},
  title={Beispiel},
  #1
}

% Lösungs-Box (für Aufgabenhefte)
\newtcolorbox{loesung}[1][]{%
  enhanced, breakable,
  colback=bfwBgAmber!40,
  colframe=bfwSuccess,
  boxrule=0.6pt,
  arc=3pt,
  left=10pt, right=10pt, top=8pt, bottom=8pt,
  fonttitle=\bfseries\color{bfwSuccess!70!black},
  title={Lösung},
  #1
}

% Glossar-Eintrag
\newcommand{\glossareintrag}[2]{%
  \par\noindent\textbf{\color{bfwPrimaryDark}#1} \enspace #2\par\smallskip
}

% Section-Stil — etwas farbiger als Standard
\usepackage{titlesec}
\titleformat{\section}
  {\Large\bfseries\color{bfwPrimaryDark}}
  {\thesection}{0.6em}{}
\titleformat{\subsection}
  {\large\bfseries\color{bfwPrimary}}
  {\thesubsection}{0.5em}{}
\titleformat{\subsubsection}
  {\normalsize\bfseries\color{bfwInk}}
  {\thesubsubsection}{0.4em}{}

% TikZ-Stil "skizzenhaft" — etwas weicher als Rough.js, aber stabil
\tikzset{
  roughbox/.style = {
    draw=bfwPrimaryDark, thick, fill=bfwBgTeal,
    rounded corners=4pt, inner sep=6pt
  },
  roughaccent/.style = {
    draw=bfwAccent, thick, fill=bfwBgAmber,
    rounded corners=4pt, inner sep=6pt
  },
  roughline/.style = {
    draw=bfwInkSoft, thick, -{Stealth[length=2.5mm]}
  }
}

% Bullet-Symbole (bewusst kein FontAwesome — vermeidet Font-Installations-Probleme)
\newcommand{\faLightbulb}{$\bullet$}
\newcommand{\faBookmark}{$\bullet$}

% Header/Footer
\usepackage{fancyhdr}
\pagestyle{fancy}
\fancyhf{}
\renewcommand{\headrulewidth}{0pt}
\fancyfoot[L]{\small\color{bfwInkSoft}\thepage}
\fancyfoot[R]{\small\color{bfwInkSoft} BFW M-064 Beschaffung im Büromanagement}


\title{\vspace*{-1.5cm}\Huge\bfseries\color{bfwPrimaryDark}Tag~15: Vertiefung, Übung \& Reflektion\\[0.4em]
  \large\color{bfwInkSoft}\textnormal{Der ganze Büroprozess auf einen Blick --- Wiederholung Tag 1--14}}
\author{\textsf{Büroprozesse gestalten und Arbeitsvorgänge organisieren \textperiodcentered{} M-067}\\
\textsf{\small\copyright\ Can Siebert\,\textperiodcentered\,2026}}
\date{17.07.2026}

\begin{document}
\maketitle
\thispagestyle{empty}

\vspace{1em}
\begin{merksatz}[title={\bfseries Worum geht es heute?}]
Der Abschlusstag bringt \emph{kein} neues Wissen --- er \emph{vernetzt} alles, was Sie in
den vergangenen 14 Tagen gelernt haben. Vom geplanten Büroraum über die gesunde,
gut organisierte Arbeit bis zur rechtssicheren Aufbewahrung und zum digitalen
Dokumentenmanagement bilden die einzelnen Themen einen \emph{durchgängigen Büroprozess}.
Dieses Handout fasst jeden Themenblock kompakt zusammen, zeigt die Querverbindungen an
einer integrierenden Fallstudie und bereitet Sie gezielt auf die IHK-Abschlussprüfung vor.
\end{merksatz}

\tableofcontents
\vspace{1em}
\hrule
\vspace{1em}

\section{Büroraumplanung \& Büroformen (Tag 1)}

Die Raumgestaltung beeinflusst \emph{Produktivität, Gesundheit und Zufriedenheit};
Leitfaden ist die \emph{DGUV Information 215-441}. Der Flächenbedarf eines Arbeitsplatzes
ergibt sich aus der Addition der \textbf{sechs Teilflächen} (Stell-, Bewegungs-,
Funktions-, Benutzer-, Verkehrswege-, Besucher-/Besprechungsfläche). Richtwerte:
Bildschirmarbeitsplatz 8--10~m\textsuperscript{2}, Großraum ca.\ 12--15~m\textsuperscript{2}.

\begin{merksatz}
Vier klassische Büroformen: \textbf{Zelle} (Ruhe, Vertraulichkeit), \textbf{Gruppe}
(Kommunikation), \textbf{Großraum} (Kosten), \textbf{Kombi} (beides). Moderne Konzepte:
reversibles Büro (flexibel umbaubar) und non-territoriales Büro/Desksharing (weniger
Plätze als Beschäftigte).
\end{merksatz}

\section{Büroausstattung \& Bildschirmarbeitsplatz (Tag 2)}

\emph{Ergonomie} heißt: die Arbeit an den Menschen anpassen. Der Bürodrehstuhl folgt
\textbf{DIN EN 1335}, der Arbeitstisch soll höhenverstellbar und reflexionsarm sein.
„So sitzen Sie richtig'': Bildschirm seitlich zum Fenster, Blickabstand 50--70~cm,
ca.\ 90°-Winkel in Ellenbogen, Hüfte und Knie. Die
\textbf{Steh-Sitz-Dynamik} lautet 60~\% sitzen, 30~\% stehen, 10~\% gehen.

\begin{eselsbruecke}
\textbf{Prüfsiegel:} \textbf{GS} = „Geprüfte Sicherheit'' (freiwillig, geprüft),
\textbf{CE} = Konformitätserklärung des Herstellers (Pflicht im EU-Binnenmarkt).
Maßgeblich ist die Arbeitsstättenverordnung (ArbStättV, Anhang Punkt~6).
\end{eselsbruecke}

\section{Arbeitsumgebung \& Umweltfaktoren (Tag 3)}

Umweltfaktoren wirken auf Gesundheit und Leistung: \textbf{Raumklima} (Lufttemperatur,
-feuchtigkeit, -bewegung; ASR~A3.5: Räume möglichst nicht über 26~°C,
ab 30~°C Maßnahmenpflicht), \textbf{Beleuchtung} (ausreichende
Lux-Werte, blend- und flimmerfrei), \textbf{Lärm} (Schalldruckpegel in dB(A)),
\textbf{Farbe} (Farbpsychologie), \textbf{Pflanzen} (Raumbegrünung) sowie Schutz vor
\textbf{Strahlung und Gefahrstoffen} (z.\,B.\ Tonerstaub).

\section{Arbeitsschutz \& Gesundheit (Tag 4)}

Drei Rechtssäulen: \textbf{ArbSchG} (Grundpflichten, \S~5 Gefährdungsbeurteilung,
\S~12 Unterweisung), \textbf{ArbZG} (Höchstarbeitszeit, Pausen, Ruhezeiten) und
\textbf{ASiG} (Betriebsarzt, Fachkraft für Arbeitssicherheit). Aufsicht: Berufs\-genossenschaft
und Gewerbeaufsichtsamt.

\begin{merksatz}
\textbf{Belastung} = von außen einwirkende Faktoren (für alle gleich); \textbf{Beanspruchung}
= individuelle Auswirkung. \textbf{Eustress} ist positiver, \textbf{Distress} negativer
(krankmachender) Stress; dauerhafter Distress kann zu \emph{Burnout} führen.
\end{merksatz}

\section{Zeitmanagement \& Aufgabenplanung (Tag 5)}

Ziel ist eine stressärmere, ausgeglichene Arbeitsweise. Zentrale Methoden:
\textbf{ALPEN} (Aufgaben, Länge, Pufferzeit, Entscheidungen, Nachkontrolle),
\textbf{Eisenhower} (wichtig\,$\times$\,dringend), \textbf{Pareto} (80/20),
\textbf{ABC-Analyse} und \textbf{Salamitaktik}. Ziele formuliert man nach \textbf{SMART}.
Wichtig: persönliche \emph{Leistungskurve} nutzen, den \emph{Sägeblatteffekt} durch eine
„stille Stunde'' vermeiden.

\section{Terminplanung \& Terminmanagement (Tag 6)}

Ein Termin besteht aus \emph{Kalenderdatum und Uhrzeit}. Man unterscheidet feste/variable,
Einzel-/Serien- und interne/externe Termine. Die \textbf{Terminüberwachung} sichert
Fristen über die \emph{Wiedervorlage}; gesetzliche Termine (Sozialabgaben, Lohnsteuer)
sind unbedingt einzuhalten. Die Urlaubsplanung richtet sich nach \S~7~BUrlG. Hilfsmittel:
Terminkalender und elektronische Tools (MS Outlook, geteilte Kalender, Doodle).

\section{Besprechungen vorbereiten \& durchführen (Tag 7)}

Über den Erfolg entscheidet die \textbf{Vorbereitung}: Datum/Uhrzeit, Raum, Teilnehmerkreis,
Zeitbedarf, \textbf{Tagesordnung (TOP)}, Unterlagen und Protokollführung klären; Einladung
fristgerecht versenden (z.\,B.\ \S~51 GmbHG für die Gesellschafterversammlung). Bei der
\textbf{Durchführung} zählen klare Rollen (Leitung, Moderation, Protokoll), Diskussions\-regeln
und der souveräne Umgang mit schwierigen Teilnehmertypen.

\section{Nachbereitung \& Protokoll (Tag 8)}

Nachbereitung heißt: Aufgaben verfolgen (To-do-/Maßnahmenliste), Protokoll erstellen und
verteilen. Ein Protokoll hat drei Funktionen: \textbf{Information, Beweis, Organisation}.

\begin{merksatz}
Fünf \textbf{Protokollarten}: Wort-, Verlaufs-, Kurz-/Ergebnis-, Beschluss- und
Gedächtnisprotokoll. Aufbau nach \textbf{DIN 5008} (Kopf, Text, Fuß). Sprache:
\emph{Präsens, sachlich, indirekte Rede mit Konjunktiv I/II}; Anträge und Beschlüsse
wörtlich.
\end{merksatz}

\section{Kommunikation \& Teamarbeit (Tag 9)}

Grundmodell ist das \textbf{Sender-Empfänger-Modell} (Störungen durch Rauschen,
Missverständnisse). Nach Schulz von Thun hat jede Nachricht \textbf{vier Seiten}
(Sach-, Selbstoffenbarungs-, Beziehungs-, Appellebene). Man unterscheidet verbale,
nonverbale und paraverbale Kommunikation. Teams durchlaufen die Phasen nach Tuckman:
\textbf{Forming, Storming, Norming, Performing, Adjourning}. Feedback gelingt mit
Ich-Botschaften, konkret und wertschätzend.

\section{Postbearbeitung: Posteingang (Tag 10)}

Der Posteingang ist der erste Umschlagplatz für Informationen (zentral/dezentral).
Arbeitsschritte: \textbf{Annehmen $\rightarrow$ Vorsortieren $\rightarrow$ Öffnen
$\rightarrow$ Kontrollieren $\rightarrow$ Stempeln $\rightarrow$ Verteilen} (dann Erfassen).
Als \emph{persönlich/vertraulich} gekennzeichnete Post wird \textbf{nicht geöffnet}.
Der Posteingangsstempel dokumentiert den Eingang; Scan/OCR ermöglicht die digitale
Archivierung.

\section{Postbearbeitung: Postausgang (Tag 11)}

Ablauf des Postausgangs: \textbf{Kontrollieren $\rightarrow$ Sortieren $\rightarrow$
Falzen $\rightarrow$ Kuvertieren $\rightarrow$ Wiegen $\rightarrow$ Frankieren
$\rightarrow$ Aufgeben}. Versandarten reichen vom Standardbrief über die
Einschreiben-Varianten bis zu KEP-Diensten. Falzarten (z.\,B.\ Wickelfalz, Zickzackfalz)
passen zu Format und Briefhülle; frankiert wird per Briefmarke, Frankiermaschine oder
Online-Frankierung. Große Mengen laufen über die \emph{Poststraße}.

\section{Dokumentenmanagement: Aufbewahrung \& Fristen (Tag 12)}

Schriftgut wird aus betrieblichen, rechtlichen und steuerlichen Gründen aufbewahrt; die
\textbf{GoBD} fordern eine ordnungsmäßige, revisionssichere (auch digitale) Aufbewahrung.
Schriftgut wird nach \emph{Wertigkeit} eingestuft (Tages-, Prüf-, Gesetzes-, Dauerwert).

\begin{merksatz}
\textbf{Aufbewahrungsfristen} (HGB \S~257, AO \S~147): \textbf{6 Jahre} für Handels-/Geschäftsbriefe,
\textbf{8 Jahre} für Buchungsbelege (neu nach BEG~IV), \textbf{10 Jahre} für Bücher,
Inventare und Jahresabschlüsse. Die Frist beginnt mit dem \emph{Schluss des Kalenderjahres},
in dem das Schriftstück entstanden ist.
\end{merksatz}

\section{Ablage \& Ablagesysteme (Tag 13)}

Eine gute Ablage macht Schriftgut schnell wiederauffindbar. \textbf{Ordnungssysteme}:
alphabetisch, numerisch, chronologisch, geografisch und alphanumerisch. \textbf{Ablagearten}:
liegende und stehende Ablage, Loseblattablage, Hängeregistratur. Eine eindeutige,
einheitliche Ordnung und ein Aktenplan sind entscheidend für die Wiederauffindbarkeit.

\section{Digitales DMS, Datenschutz \& Umweltschutz (Tag 14)}

Ein \textbf{Dokumentenmanagementsystem (DMS)} ermöglicht schnellen Zugriff, Volltextsuche,
Mehrfachnutzung, Workflows und revisionssichere Archivierung. Der \textbf{Datenschutz}
(DSGVO, BDSG) verlangt den Schutz personenbezogener Daten durch \emph{technisch-organisatorische
Maßnahmen (TOM)} wie Zugriffsrechte, Verschlüsselung und Löschkonzepte. Der
\textbf{Umweltschutz} setzt auf das papierlose Büro, Recyclingpapier, Mülltrennung und
Energie sparen --- erkennbar z.\,B.\ am Umweltsiegel \emph{Blauer Engel}.

\section{Integrierende Fallstudie: ein Arbeitstag, viele Themen}

\begin{beispiel}
Bei der \emph{Lindner \& Partner GmbH} (Steuerberatung, 18 Beschäftigte) erleben Sie als
Kaufmann/-frau für Büromanagement einen typischen Tag --- und brauchen fast das ganze
Modul:
\begin{itemize}
  \item \textbf{Raum \& Ergonomie} (Tag 1--3): Für eine neue Kollegin empfehlen Sie ein
    \emph{Kombibüro} und korrigieren Bildschirmstellung, Stuhl und Raumtemperatur.
  \item \textbf{Zeit \& Termine} (Tag 5--6): Sie priorisieren mit dem \emph{Eisenhower-Prinzip}
    und bereiten eine Besprechung samt \emph{Tagesordnung} vor (Tag 7).
  \item \textbf{Protokoll \& Kommunikation} (Tag 8--9): Sie führen ein \emph{Beschlussprotokoll}
    (Präsens, indirekte Rede) und entschärfen einen Konflikt mit dem Modell der
    \emph{vier Seiten einer Nachricht}.
  \item \textbf{Post, Fristen \& DMS} (Tag 10--14): Ein \emph{vertrauliches Einschreiben}
    leiten Sie ungeöffnet weiter, eine Eingangsrechnung legen Sie ab (Aufbewahrung
    \emph{8 Jahre}, Fristbeginn Jahresende) und überführen die Akten datenschutz\-gerecht
    in ein digitales DMS.
\end{itemize}
So zeigt sich: Die Themen des Moduls greifen im Alltag \emph{ineinander} --- genau das
prüft auch die IHK.
\end{beispiel}

\newpage

\section{Ausblick: die IHK-Abschlussprüfung}

Mit diesem Tag endet das Modul. In der schriftlichen IHK-Prüfung kommt es darauf an, das
Wissen \emph{anzuwenden} und Antworten \emph{zu begründen} --- bloße Stichworte reichen
selten. Achten Sie genau auf die \textbf{Operatoren}:

\begin{itemize}
  \item \textbf{nennen/aufzählen} --- nur die geforderten Begriffe auflisten.
  \item \textbf{erläutern/beschreiben} --- erklären und begründen, in ganzen Sätzen.
  \item \textbf{beurteilen/bewerten} --- eine eigene, mit Argumenten gestützte Stellungnahme.
\end{itemize}

\begin{eselsbruecke}
\textbf{Prüfungstipps:} Zuerst die ganze Aufgabe lesen, Punktzahl als Hinweis auf den
Umfang nutzen, Zeit nach Punkten einteilen, leichte Aufgaben zuerst, am Ende auf
Vollständigkeit prüfen. Der IHK-Notenschlüssel beginnt bei \emph{ausreichend} ab 50~\%.
\end{eselsbruecke}

Nutzen Sie das Abschlussquiz (30 Fragen über alle 14 Themen) und das Aufgabenheft, um
Ihre Wissenslücken gezielt zu schließen. \textbf{Herzlichen Glückwunsch zum Kursabschluss
--- und viel Erfolg in der Prüfung!}

\newpage

\section*{Glossar --- zentrale Begriffe des Moduls}
\addcontentsline{toc}{section}{Glossar}

\glossareintrag{Teilflächen}{Die sechs Bestandteile (Stell-, Bewegungs-, Funktions-, Benutzer-, Verkehrswege-, Besucher-/Besprechungsfläche), aus denen sich der Flächenbedarf ergibt.}

\glossareintrag{Kombibüro}{Einzelbüros an der Fensterfront plus zentrale Multifunktionszone; vereint Einzel- und Teamarbeit.}

\glossareintrag{Ergonomie}{Anpassung der Arbeit an den Menschen; Bürodrehstuhl nach DIN EN 1335.}

\glossareintrag{Steh-Sitz-Dynamik}{Wechsel der Körperhaltung im Verhältnis ca.\ 60~\% sitzen, 30~\% stehen, 10~\% gehen.}

\glossareintrag{ASR A3.5}{Technische Regel zum Raumklima; Räume möglichst nicht über 26~°C, ab 30~°C Maßnahmenpflicht.}

\glossareintrag{Gefährdungsbeurteilung}{Pflicht des Arbeitgebers nach \S~5 ArbSchG, Gefahren am Arbeitsplatz systematisch zu ermitteln und zu bewerten.}

\glossareintrag{Belastung / Beanspruchung}{Belastung = äußere Einwirkung (für alle gleich); Beanspruchung = individuelle Auswirkung auf die Person.}

\glossareintrag{ALPEN-Methode}{Tagesplanung: Aufgaben notieren, Länge schätzen, Pufferzeiten, Entscheidungen/Prioritäten, Nachkontrolle.}

\glossareintrag{Eisenhower-Prinzip}{Priorisierung von Aufgaben nach Wichtigkeit und Dringlichkeit (A/B/C/Papierkorb).}

\glossareintrag{SMART}{Regel zur Zielformulierung: spezifisch, messbar, attraktiv, realistisch, terminiert.}

\glossareintrag{Wiedervorlage}{Verfahren der Terminüberwachung; Vorgänge werden zum festgelegten Termin erneut vorgelegt.}

\glossareintrag{Tagesordnung (TOP)}{Geordnete Liste der Besprechungspunkte; Teil von Einladung und Vorbereitung.}

\glossareintrag{Protokoll}{Schriftliche Aufzeichnung einer Besprechung mit den Funktionen Information, Beweis und Organisation.}

\glossareintrag{Beschlussprotokoll}{Protokollart, die ausschließlich die gefassten Beschlüsse festhält.}

\glossareintrag{Vier Seiten einer Nachricht}{Modell von Schulz von Thun: Sach-, Selbstoffenbarungs-, Beziehungs- und Appellebene.}

\glossareintrag{Tuckman-Phasen}{Teamentwicklung: Forming, Storming, Norming, Performing, Adjourning.}

\glossareintrag{Posteingangsstempel}{Stempel zur Dokumentation von Eingangsdatum (und Bearbeitung) eines Schriftstücks.}

\glossareintrag{Poststraße}{Maschineller Ablauf des Postausgangs (Falzen, Kuvertieren, Verschließen, Frankieren).}

\glossareintrag{GoBD}{Grundsätze ordnungsmäßiger Buchführung und Aufbewahrung in digitaler Form (Unveränderbarkeit, Revisionssicherheit).}

\glossareintrag{Aufbewahrungsfristen}{HGB \S~257 / AO \S~147: 6, 8 oder 10 Jahre; Fristbeginn mit Schluss des Kalenderjahres.}

\glossareintrag{Ordnungssysteme}{Verfahren der Ablage: alphabetisch, numerisch, chronologisch, geografisch, alphanumerisch.}

\glossareintrag{DMS}{Dokumentenmanagementsystem zur digitalen Erfassung, Verwaltung und Archivierung von Dokumenten.}

\glossareintrag{TOM}{Technisch-organisatorische Maßnahmen zum Schutz personenbezogener Daten nach DSGVO.}

\glossareintrag{Blauer Engel}{Bekanntes Umweltsiegel für umweltfreundliche Produkte (z.\,B.\ Recyclingpapier).}

\end{document}
